\documentclass[12pt,reqno]{amsart}
\usepackage{math-review}
% Project-wide mathematical notation. No claims are certified here.

\title{The Schur Complement -- Minimal}
\author{Known-mathematics exposition benchmark}
\date{}
\begin{document}
\maketitle
\section{Main Results}
The cross term in a quadratic form can be eliminated without losing the
nonzero-vector condition. The resulting criterion is classical.
\begin{theorem}[Schur-complement criterion]\label{thm:main}
Let $p,q\geq1$ be integers. Let $A\in\mathbb R^{p\times p}$ and
$C\in\mathbb R^{q\times q}$ be symmetric, and let $B\in\mathbb R^{p\times q}$.
Assume that $x^{\mathsf T}Ax>0$ for every $0\ne x\in\mathbb R^p$. Define
\[
 M=\begin{pmatrix}A&B\\ B^{\mathsf T}&C\end{pmatrix},
 \qquad S=C-B^{\mathsf T}A^{-1}B.
\]
Then $M$ is positive definite if and only if $S$ is positive definite.
Here a real symmetric matrix $H$ is positive definite when
$v^{\mathsf T}Hv>0$ for every nonzero real vector $v$ of the appropriate size.
\end{theorem}
The positivity of $A$ is essential. Invertibility alone allows the square
identity below, but does not control the sign of its first term.

\section{Key Mathematical Ingredients}
The reader may test positivity by a quadratic form. What needs refreshing
is the choice that removes the mixed term, rather than the name of a matrix.
\begin{lemma}[Square completion]\label{lem:square}
Let $A\in\mathbb R^{p\times p}$ be symmetric and invertible,
$B\in\mathbb R^{p\times q}$ and $C\in\mathbb R^{q\times q}$ symmetric.
For all $x\in\mathbb R^p$ and $y\in\mathbb R^q$, putting
$z=x+A^{-1}By$ gives
\[
 x^{\mathsf T}Ax+2x^{\mathsf T}By+y^{\mathsf T}Cy
 =z^{\mathsf T}Az+y^{\mathsf T}(C-B^{\mathsf T}A^{-1}B)y.
\]
\end{lemma}
Denote the left side of this identity by $Q(x,y)$.

Its output is a sum of two uncoupled forms. The next ingredient explains
why this sum still tests every nonzero vector.
\begin{lemma}\label{lem:inverse}
For every real matrix $K\in\mathbb R^{p\times q}$, the linear map
\[
 (x,y)\longmapsto (x+Ky,y)
\]
is invertible, with inverse $(z,y)\mapsto(z-Ky,y)$. Consequently it takes
nonzero vectors to nonzero vectors, in both directions.
\end{lemma}
Here we use $K=A^{-1}B$. The inverse map is the bridge between the
original positivity test and the two remaining quadratic forms.

\section{Connecting the Ingredients}
For the first direction, $A>0$ and $S>0$ give, for every $(x,y)\ne0$,
\[
 Q(x,y)=z^{\mathsf T}Az+y^{\mathsf T}Sy>0.
\]
Indeed $(z,y)\ne0$ by Lemma~\ref{lem:inverse}, so at least one positive
term is nonzero. Both positivity assumptions are used jointly.

For the reverse direction, fix an arbitrary $y\ne0$ and choose
\[
 x=-A^{-1}By,\qquad z=0.
\]
This is the decisive test: it removes the already controlled square.
The original pair is nonzero, and $M>0$ therefore yields
\[
 0<Q(-A^{-1}By,y)=y^{\mathsf T}Sy.
\]
Since this holds for every nonzero $y$, it proves $S>0$.

For a boundary check, take $A=-1$, $B=0$, $C=1$. Then $S=1>0$, but
$M=\operatorname{diag}(-1,1)$ is not positive definite. This explains
why the theorem does not replace $A>0$ by mere invertibility.

\nocite{S-guide}
\bibliographystyle{plain}
\bibliography{../generated/references}
\end{document}
